\documentclass[12pt, a4paper]{scrbook}
\usepackage[utf8]{inputenc}
\usepackage{csquotes}
\usepackage[german]{babel}
\usepackage{hyperref}
\usepackage[onehalfspacing]{setspace}
\usepackage{geometry}
\usepackage{color}
\usepackage{listings}
\usepackage{graphicx}
\usepackage{acronym}
\usepackage[backend=biber,
bibstyle=numeric, 
citestyle=authortitle,
natbib=true, 
hyperref=true,
]{biblatex}
\addbibresource{lit.bib} 
\setcounter{secnumdepth}{4}
\setcounter{tocdepth}{4}
\geometry{
left=2.5cm,
right=2.5cm,
top=2.5cm,
bottom=2.5cm,
bindingoffset=5mm,
}
\pagestyle{empty}
\begin{document}
%\documentclass[titlepage, 12pt]{scrbook}
%\usepackage[utf8]{inputenc}
%\usepackage[german]{babel}
%\usepackage{uarial}
%\usepackage{color}
%\renewcommand{\familydefault}{\sfdefault}
%\usepackage{fancyhdr}
%\usepackage{graphicx}
%\pagestyle{empty}
%\begin{document}
\begin{titlepage}
	\begin{flushleft}
	\begin{figure}
		\hspace*{-0,5cm}
		\includegraphics[scale=0.25]{Bilder/DHBW_logo.jpg} \hspace*{5cm}
		\includegraphics[scale=0.25]{Bilder/Atos_logo.png}
	\end{figure}
	\end{flushleft}
	\vspace*{-0.6cm}
	\begin{center}
	\textcolor{red}{Titel der großen Studienarbeit} \par \vspace*{0,5cm}
	Projektarbeit \par \vspace*{2cm}
	des Studienganges \textcolor{red}{Angewandte Informatik / Betriebliches Informationsmanagement}
	an der Dualen Hochschule Baden-Württemberg Mannheim \par \vspace*{1cm}
	von \par \vspace*{0,5cm}
	\textcolor{red}{Maximilian Ludwig, Kevin Wrona, Fabian Brandmüller} \par \vspace*{1cm}
	\today \par \vspace*{2cm}
	\begin{tabular}{l@{\hspace{3cm}}r}
		Bearbeitungszeitraum & \textcolor{red}{23.09.2019 - 20.04.2019} \\
		Betreuer der DHBW & \textcolor{red}{Eckhard Kruse} \\[1cm] 
	\end{tabular}
	\end{center}
\end{titlepage}
%\end{document}
\setlength{\parindent}{0em} 
\renewcommand\thechapter{\Roman{chapter}}
\pagenumbering{Roman}
\let\cleardoublepage\relax
\section*{Erklärung}
Ich versichere hiermit, dass ich meine Projektarbeit mit dem Thema: ''Titel der großen Studienarbeit'' selbstständig verfasst und keine anderen als die angegeben Quellen und Hilfsmittel benutzt habe. 
\newline
\newline
\newline
\newline
---------------------------------------------       ------------------------------------------ \newline
Ort	\hspace{2cm}		Datum\hspace{3,5 cm}				    Unterschrift
\newpage
\section*{Abstract}


\newpage
\begingroup
\renewcommand*{\chapterpagestyle}{empty}
\pagestyle{empty}
\tableofcontents
\listoffigures
\section*{Abkürzungen}
\begin{acronym}[Bash]
% \acro{VCS}{Version Control System}
\end{acronym}
\endgroup
\newpage
\pagestyle{plain}
\setcounter{page}{1}
\chapter{Einleitung}
''Several researchers have stated that facial expression recognition appears to play one of the most important roles in human communication'' 
\footcite[Vgl.][1]{FaceRec}
Dieses Zitat von Katherine B. Leeland gibt einen Einblick in die Relevanz der Emotionserkennung für den Menschen. Dabei ist diese Relevanz nicht erst in jüngerer zeit entstanden. Bereits Darwin stellte die Frage, ob von den Gesichtsausdrücken einer Person nicht auch der Emotionale Zustand abgeleitet werden kann.
\footcite[Vgl.][2]{FaceRec}
Einen solchen Zustand von einem Mitmenschen abzulesen ist jedoch nicht trivial. Dies liegt auch daran, dass die Übergänge von verschiedenen Emotionen durch leichte Veränderungen der Mimik ausgedrückt werden. Zum Beispiel indem eine Person die Lippen zusammen presst und die Augen zusammen kneift bei Wut, oder die Mundwinkel nach unten gezogen werden bei Trauer. 
\footcite[Vgl.][249]{HandbookFaceRec}
Durch solche kleinen Änderungen kann ein sonst neutral wirkendes Gesicht zu einem wütenden oder traurigen werden.
Emotionserkennungssoftware gibt es bereits und wird auch in der Wirtschaft eingesetzt. Die Anwendungsgebiete reichen dabei von Jobinterviews, in denen analysiert wird in wie weit die Bewerber zu dem jeweiligen Job passen 
\footcite[Vgl.][]{mixedArticle}
bis hin zur Automobilindustrie. Dort wird mittels geegineter Sensorik versucht die Emotion und somit der physiologische Zustand des Autofahrers zu analysieren.
\footcite[Vgl.][Herausforderung]{Frauenhofer}
Mit diesen Daten können dann wiederum Funktionen eingebaut werden wie Warnsysteme, die den Fahrer daraufhinweisen, dass sein Zustand ungeeignet zum Betrieb eines Kraftfahrzeugs ist. Jedoch werden solche Einsatzszenarien auch durchaus kontrovers diskutiert. Auf Kritik stößt unter anderem dass die sogenannten ''Basisemotionen'', die verwendet werden um den KIs Emotionserkennung beizubringen, selber umstritten sind. 
\footcite[Vgl.][]{SZ}
Aber auch ethische bedenken werden zunehmend geäußert, vor allem bezüglich der Anwendungsgebiete. Denn je nach Emotion die erkannt werden soll, liegt die Fehlerrate sehr hoch. So hat das Frauenhofer institut, welches an Einsatzgebieten von Emotionserkennung in Fahrezugen arbeitet, festgestellt, dass eine Emotionserkennung je nach Zielemotion eine Vorhersagekraft zwischen 6 und 95\% haben kann.
\footcite[Vgl.][Ergebnis]{Frauenhofer}
Diese negativen Aspekte treffen jedoch nur teilweise auf das hier behandelte Forschungsprojekt zu, wie im Folgenden dargelegt werden soll:
Zwar werden in dieser Arbeit wie im folgenden Kapitel beschrieben werden wird die Basisemotionen als Referenz herangezogen, da es jedoch momentan keine praktikable Möglichkeit gibt mit einem ähnlichen Workload eine solche Lösung zu erstellen, sind diese unumgänglich. Auf der anderen Seite ist diese Studienarbeit kein wissenschaftliches Werk über die psychologischen Emotionen die hinter verschiedenen Gesichtsausdrücken stehen, auch wenn sie dies indirekt thematisiert. Es geht hier darum die These zu verifizieren oder zu widerlegen, ob es möglich ist mit einer künstlichen Intelligenz zu erkennen, in wie weit ein Pokerface einer Person anhand dessen Gesichtsausdruck abgelesen werden kann. Diese Aufgabe hat also nicht den Anspruch eine Emotion korrekt vorherzusagen, sondern vorherzusagen wann keine Emotion vorliegt. Denn Ein Pokerface wird im Allgemeinen als eine emotionsloser Gesichtsausdruck definiert. Dies impliziert, dass eine Person mit Pokerface keine Emotion zu erkennen gibt. Daher könnte bei einer nicht messbaren, oder neutralen Emotion ein Pokerface vorliegen. Dies mittels KI zu testen und einen sogenannten ''Pokerface Detektor'' zu entwicklen ist daher Ziel dieser Arbeit. Mit diesem Pokerface Detektor sind verschiedenste Einsatzmöglichkeiten in der Praxis denkbar. Diese haben eine gewisse Schnittmenge mit denen von ''normaler'' Emotionserkennungs software, jedoch gibt es auch einige weitere. Im Folgenden werden einige mögliche Szenarien expliziert:
%Beispiele ausführen
\begin{itemize}
	\item{Polizeiverhöre}
\end{itemize}
Es ist denkbar, dass ein Detektor wie der, der in dieser Arbeit entwickelt werden wird bei Polizeiverhören eingesetzt wird. Für die Beamten kann es nicht immer direkt ersichtlich sein, ob der verhörte mit einem Pokerface lügt, oder die Wahrheit sagt. Da die Möglichkeiten dies zu prüfen ebenfalls nicht zahlreich oder einfach verfügbar sind, wäre es eine Vereinfachung für die Beamten wenn eine einfache Webcam zusammen mit einem Computer reichen würde um einen Lügner zu entlarven.
\begin{itemize}
	\item{Gerichtsverhandlungen}
\end{itemize}
Das zweite Einsatzgebiet ist ähnlich zu dem ersten. Bei Gerichtsverhandlungen gelten die gleichen Vorraussetzungen wie bei einem Verhöh der Polizei. Zwar müssen die hier vorgeladenen eine eidestattliche Erklärung abgeben nur die Wahrheit zu sagen, jedoch ist zu bezweifeln ob dies auch jeder so handhabt. 
Nun soll nicht der Eindruck entstehen dass das hier gebaute Werkzeug ein Lügendetektor ist. Es ist ebenfalls nicht möglich, dass von einem Pokerface immer auf eine Lüge geschlossen werden kann. Jedoch ist ein Pokerface ein Zeichen dafür, dass sich diese Person ihren emotionalen Zustand nicht anmerken lassen möchte. Und dies wiederum deutet eher daraufhin dass die Person nicht die Wahrheit sagt oder nur teilweise. 
Abgesehen von Einsatzgebieten die zur Entlarvung von Lügen führen kann auch ein klassicheres Szenario verwendet werden:
\begin{itemize}
	\item{Pokerspiel}
\end{itemize}
Es ist anzunehmen, dass der erste Begriff der mit dem Wort Pokerface in Verbindung gebracht wird, das Pokerspiel selber ist. Und auch in diesem kann ein Pokerface Detektor nützlich sein. So kann ein Mitspieler zum Beispiel mittels einer Kamera das Gesicht des Gegenübers scannen und analysieren ob ein Pokerface vorliegt oder nicht, und dementsprechend agieren.
Nachdem verschiedene Anwendungsszenarien beleuchtet wurden wird im Folgenden die konkrete Forschungsfrage beleuchtet.
%Irgendwas vergessen?
\section{Aufgabenstellung}
 Das Projekt selber wird an der DHBW in Mannheim durchgeführt und betreut von Prof. Dr. Erckhard Kruse.
 Wie eingangs erwähnt soll mittels künstlicher Intelligenz ein Pokerface erkannt werden. 
 Dafür soll wiederum eine Bilderkennungssoftware angefertigt werden, die Bilder oder Videos als Eingae von einer Webcam oder anderen Kamera bekommt, und das Bild analysiert. Je nachdem welche emotionen gezeigt werden wird dann von der Software ein Rückschluss gezogen auf ein eventuell vorhandens Pokerface. Ein konkretes Einsatzgfebiert nach Abschluss der Entwickliung ist nicht vorgesehen, da es sich um ein Forschungsprojekt handelt. jedoch sind wie bereits beschrieben einige verschiedene Einsatzmöglichkeiten denkbar, an die das Werkzeug leicht angepasst werden kann. 

%\section{DevOps}
%\begin{figure}
%\includegraphics[width=\linewidth]{Bilder/DevOps-LifeCycle-Features-Naresh-IT.png}
%\caption{ DevOps Lifecyle \newline Quelle: https://nareshit.com/wp-content/uploads/2018/09/DevOps-LifeCycle-Features-Naresh-IT.png }
%\label{fig:Devops-Lifecycle}
%\end{figure}


\let\cleardoublepage\relax
\chapter{Theorie}
Technologien todo:
\begin{itemize}

\item Poker erklären (Wenn wir Poker als testusecase benutzen)
\end{itemize}

\section{Stand der Technik}
In diesem Abschnitt soll der aktuelle Forschnungs und Entwicklungsstand thematisiert werden im Bereich Emotionserkennung. Jedoch muss dafür erst einmal eine Unterscheidung der Begrifflichkeiten Emotionserkennung und Geischtskernnung erfolgen, da beide gebiete Überschneidungen haben, jedoch inhaltlich und von ihren Zielen verschieden sind.
\subsection{Emotionserkennung}

usecase für emotion recognition

\subsection{Geischtserkennung}
Gesichtserkennung ist eine Disziplin der Informatik in der es darum geht Gesichter wieder zu erkennen, und gegebenanfalls verschiedenen Personen zuzuordnen. Dabei lässt sich der Prozess der Gesichtserkennung in vier Phasen einteilen, Face ''detection'', ''alignment'', ''feature extraction'' und ''matching''.
\footcite[Vgl. ][2]{HandbookFaceRec}
Die Detection Phase ist dafür verantwortlich um zu erkennen ob Gesichter vorhanden sind in einem Bild, oder aber Video.
\footcite[Vgl. ][2]{HandbookFaceRec}
 In der darauffolgenden Alignment Phase hingegen wird die Lokalisierung der Geischter genauer, indem geischtskomponenten wie Augen, Augenbrauen, oder die Nase genauer lokalisiert werden. Dabei wird das Bild oder Video ebenfalls  normalisiert, indem z.B. die Bildbeleuchtung angepasst wird.
 \footcite[Vgl. ][2]{HandbookFaceRec}
 In der Feature extraction hingegen werden die verschiedenen Gesichtskomponenten wie Augen, Nase, Mund, dem Bild oder Video entnommen. Dies ist ein enorm wichtiger Schritt für weitere Prozesse wie Eye Tracking oder Face Tracking. Alternativ kann sogar eine bestimmte Person anhand der extraierten Merkmale erkannt werden.
\footcite[Vgl. ][Abstract]{IEEE}
 %https://ieeexplore.ieee.org/document/7887965
 In der letzten Phase, dem Matching, geht es darum die gewonnen Daten mit den in der Datenbank vorhandenen Gesichtern abzugleichen. Wenn eine genügende Übereinstimmung gefunden wurde, wird ein Match mit einer Person ausgegeben.
  \footcite[Vgl. ][3]{HandbookFaceRec}
 
 
\subsection{Gemeinsamkeiten und Unterschiede}
\subsection{Weitere Projekte}
\section{Emotionen}
In diesem Unterakptiel nun sollen Emotionen an sich thematisiert werden, da diese maßgeblich sind für das zu entwickelnde Tool. Grundsätzlich gibt es viele verschieden Ansätze Emotionen zu definieren, bzw. Emotionen einzuteilen. Eine gängige Einteilung ist dabei die verschiedenen Emotionen in acht Bereiche einzuteilen. Diese Einteilung wurden 1984 von Plutchik postuliert und beinhaltet die Emotionskategorien Angst, Wut, Freude, Trauer, Akzeptanz, Ekel, Erwartung und Überraschung.
\footcite[Vgl. ][3]{FaceRec}
Jedoch ist dies bei weitem nicht die einzige mögliche Einteilung. Als weiteres Beispiel teilte MacLean  die Emotionen in lediglich sechs Kategorien ein, Verlangen, Wut, Angst, Niedergeschlagenheit, Freude und Zuneigung.
\footcite[Vgl. ][3]{FaceRec}
Wie sich bereits an den beiden Beispielen zeigt, geht die Meinungen der Forscher  dabei sehr stark auseinander, welche und wie viele Emotionen zu den sogenannten ''Basis Emotionen'' gehören. In dieser Arbeit werden die Emotionen in sechs Kategorien eingeteilt, in Wut, Trauer, Freude, Ekel, Überraschung und Neutral. Diese Einteilung entspricht an sich keiner von uns gefundenen gängigen Einteilung, jedoch wurde diese aus den folgenden Gründen gewählt:
Die hier genannten Emotionen lassen sich gut anhand von Bildern erlernen, da diese zum Teil komplementär und somit eindeutig sind. Beispielsweise ist es einfacher ein freudiges Gesicht zu erkennen, oder ein überraschtes, als ein Gesicht mit dem Emotionalen Ausdruck Akzeptanz. Des Weiteren wurde der Ausdruck ''Neutral'' hinzugefügt. Neutral rerpäsentiert ein emotionsloses Gesicht, bzw. ein Pokerface. Außerdem sind dieses Arten der Gesichtsausdrücke auch nicht selten anzutreffen bei dem Test Usecase dieser Arbeit,

%dem Texas Holdem Poker.
\section{Python}
\begin{itemize}
\item Python (Als ganzes warum Python, Libraries erklären, warum python verwendet wurde und state oif the art ist
\item Tesnorflow und Kreas wenn wir es verwednern aber auch wenn wir es nicht machen
\item verwendeter Datensatz
\item CNN --> Wie arbeitet die KI bei Gesichtserkennung ?
\end{itemize}
\let\cleardoublepage\relax
\newpage
\chapter{Methode}
Diese Arbeit soll methodisch mit der MoSCoW Priorisierung bearbeitet werden. Diese Art der Priorisierung teilt die zu bearbeitenden Anforderungen in vier Kategorien ein:
\footcite[vgl.][90]{Projektmanagement}
\begin{itemize}
\item Must - Core Anforderungen die unbedingt umgesetzt werden müssen
\item Should - Anforderungen die ebenfalls umgesetzt werden müssen, jedoch um Nachhinein noch durch Change Request verändert werden können.
\item Could - Anforderungen die Nach den Must und Should Anforderungen umgesetzt werden sollen, sofern noch Ressourcen und Zeit vorhanden sind um diese zu bearbeiten
\item Won't - Anforderungen die nicht in diesem Projekt bzw. Release erfolgen, jedoch in einer zukünftigen Version bearbeitet werden sollen. 
\end{itemize}


\begin{itemize}
\item Must
\begin{itemize}
\item placeholder
\end{itemize}
\item Should
\begin{itemize}
\item placeholder
\end{itemize}
\item Could
\begin{itemize}
\item placeholder
\end{itemize}
\item Won't
\begin{itemize}
\item placeholder
\end{itemize}
\end{itemize}
Dabei sollen die einzelnen Anforderungen entsprechend ihrer Priorität abgearbeitet werden. So kann am Ende der Erfolg der Arbeit deutlich besser eingeordnet werden

\let\cleardoublepage\relax
\chapter{Ergebnis}
%\section{test1}
%\subsection{test2}
%\subsubsection{test3}
%\paragraph*{Gitea Setup}

\let\cleardoublepage\relax
\chapter{Diskussion}
Das nunmehr letzte Kapitel soll sich mit der kurzen Zusammenfassung der Ergebnisse des letztens Teils und deren Bewertung widmen. Des Weiteren sollen die angewandten Methoden reflektiert werden, offene Fragen beantwortet und auch weitere Punkte aufgezeigt werden die verbessert oder noch implementiert werden können. Dazu soll zunächst die Ergebnisse kurz zusammengefasst werden.
\section{Reflexion der Ergebnisse}
\section{Reflexion Vorgehen}
Mehr darauf eingehen dass das Kontrovers ist und auch die basisemotionen kopntrovers sind --  aber keine andere Möglichkeit vorhanden 
\section{Reflexion der Literatur}
Bezüglich der Literatur ergeben sich nun einige Schwierigkeiten. Dies liegt unter anderem daran, dass das generelle Thema der Gesichts und Emotionserkennung immer noch vor allem aus psychologischer Sicht in der Literatur behandelt wurde. Zwar gibt es Fachbücher auch aus informationstechnischer Sicht, welceh ebenfalls in dieser Arbeit verwedet wurden. 
\section{Offene Implikationen}


\let\cleardoublepage\relax
\pagestyle{empty}
\newpage
\pagestyle{empty}
\printbibheading
\printbibliography[type=book,heading=subbibliography,title={Literaturquellen}]
\printbibliography[type=misc,heading=subbibliography,title={Online Quellen}]
\pagestyle{empty}
\newpage
\pagestyle{empty}

\end{document}
